\chapter{导论(Introduction)}
\intro{凡筒井皆用机械，利之所在，人无不知。
}\rightline{——苏轼}

\section{研究对象}
{\Huge{Machinary(机械)}}————Machine(机器) and Mechanism(机构)

\subsection{机器的三个特征}
\begin{enumerate}
    \item 人为实物的组合
    \item 各组成部分间有确定的相互运动
    \item 能做有用功，完成信息的传递和能量的转换
\end{enumerate}

如果只满足前两个特征，就是机构，如果满足三个特征，就是机器。

\subsection{现代机械的组成}
\begin{enumerate}
    \item 原动机（Prime mover）
    \item 执行机构(Working system)
    \item 传动机构(Transmission system)
    \item 控制机构(Control system)
\end{enumerate}

\section{相关概念}
\subsection{构件}
\begin{definition}
    构件是一个机器的运动单元体，具有确定的运动。
\end{definition}
\subsubsection{构件的分类}
\begin{enumerate}
    \item 机架：固定不动的构件
    \item 原动件：驱动力（矩）作用的构件
    \item 从动件：随原动件运动而运动的构件
\end{enumerate}

\subsection{零件}
\begin{definition}
    零件是一个机器的制造单元体，不可拆分。
\end{definition}

\subsection{部件}
\begin{definition}
    为完成特定的功能在结构上组合在一起并协同工作的零件组合。
\end{definition}
